\section{Lema de Burnside}

\textbf{Problema base.} Contar quantas configurações são \textbf{diferentes} quando algumas são consideradas
iguais por simetria. Ex.: colares que viram o mesmo colar quando giramos.

\textbf{Ideia.} Para cada simetria, conte quantas configurações continuam \textbf{iguais} depois de aplicá-la.
A resposta é a média dessas contagens (cada configuração diferente é contada exatamente $|G|$ vezes na soma):
\[
\boxed{\text{resposta} = \frac{1}{|G|} \sum_{g \in G} \mathrm{fixas}(g)}
\]
\begin{itemize}
  \item $G$: o conjunto de simetrias, incluindo ``não fazer nada'' (ex.: girar $0, 1, \dots, n-1$ posições); $|G|$ é quantas são.
  \item $g$: uma simetria de $G$.
  \item $\mathrm{fixas}(g)$: quantas configurações não mudam quando aplicamos $g$.
\end{itemize}

\textbf{Como contar $\mathrm{fixas}(g)$ em colorações.} Siga para onde $g$ leva cada posição: as posições
se dividem em ciclos. Para a coloração não mudar, cada ciclo precisa ter uma cor só. Com $c(g)$ ciclos e $k$ cores:
\[ \mathrm{fixas}(g) = k^{c(g)}. \]

\textbf{Exemplo.} Colares de $4$ contas com $2$ cores ($|G| = 4$ rotações):
\begin{itemize}
  \item girar $0$: cada posição é um ciclo, $4$ ciclos: $2^4 = 16$;
  \item girar $1$: $0 \to 1 \to 2 \to 3 \to 0$, $1$ ciclo: $2^1 = 2$;
  \item girar $2$: $(0\,2)(1\,3)$, $2$ ciclos: $2^2 = 4$;
  \item girar $3$: $1$ ciclo: $2^1 = 2$.
\end{itemize}
Resposta: $(16 + 2 + 4 + 2)/4 = 6$, que são $0000, 0001, 0011, 0101, 0111, 1111$.

\subsection*{Casos clássicos ($n$ posições, $k$ cores)}

\begin{itemize}
  \item \textbf{Colar} (só girar, $|G| = n$): girar $i$ posições forma $\gcd(i, n)$ ciclos.
        Ex.: $n = 6$, girar $2$: $(0\,2\,4)(1\,3\,5)$, e $\gcd(2, 6) = 2$.
        \[ \frac{1}{n} \sum_{i=0}^{n-1} k^{\gcd(i, n)} = \frac{1}{n} \sum_{d \mid n} \varphi\!\left(\tfrac{n}{d}\right) k^{d}. \]
        A segunda forma agrupa as rotações pelo número de ciclos $d$: exatamente $\varphi(n/d)$ valores de $i$ têm
        $\gcd(i, n) = d$. Ela só passa pelos divisores de $n$, útil para $n$ grande.

  \item \textbf{Pulseira} (pode girar e virar, $|G| = 2n$): são as $n$ rotações do colar mais $n$ reflexões. Uma reflexão
        espelha num eixo: as contas em cima do eixo ficam paradas (ciclos de $1$) e as outras trocam em pares (ciclos de $2$).
        Com $n$ ímpar, cada eixo passa por $1$ conta: $\frac{n+1}{2}$ ciclos. Com $n$ par, metade dos eixos passa por
        $2$ contas ($\frac{n}{2} + 1$ ciclos) e metade passa entre contas ($\frac{n}{2}$ ciclos). Somando as reflexões em $R$:
        \[
        \frac{1}{2n}\left(\sum_{i=0}^{n-1} k^{\gcd(i, n)} + R\right), \qquad
        R = \begin{cases}
          n\,k^{(n+1)/2}, & n \text{ ímpar} \\
          \frac{n}{2}\left(k^{n/2+1} + k^{n/2}\right), & n \text{ par}
        \end{cases}
        \]

  \item \textbf{Grade $n \times n$} (girar $0^\circ, 90^\circ, 180^\circ, 270^\circ$, $|G| = 4$): em $0^\circ$ cada célula é um
        ciclo ($n^2$); em $90^\circ$ e $270^\circ$ as células andam em grupos de $4$ ($\lceil n^2/4 \rceil$ ciclos); em
        $180^\circ$ trocam em pares ($\lceil n^2/2 \rceil$ ciclos). O $\lceil \cdot \rceil$ conta o centro, que fica parado
        quando $n$ é ímpar.
        \[ \frac{1}{4}\left(k^{n^2} + 2\,k^{\lceil n^2/4 \rceil} + k^{\lceil n^2/2 \rceil}\right). \]
\end{itemize}
